% !TeX program = pdflatex  (compiler avec -shell-escape pour minted)
\documentclass[aspectratio=169,10pt]{beamer}

% ============================================================
%  THEME
% ============================================================
\usetheme[progressbar=frametitle,
          numbering=fraction,
          block=fill,
          sectionpage=progressbar]{metropolis}

% ------------------------------------------------------------
%  PALETTE : une seule couleur d'accent + niveaux de gris.
%  Pour changer toute l'identite du deck, modifier ces 2 lignes.
% ------------------------------------------------------------
\definecolor{accent}{HTML}{A6192E}   % crimson sobre
\definecolor{ink}{HTML}{1A1A1A}      % noir du texte
\definecolor{ash}{HTML}{6E6E6E}      % gris secondaire
\definecolor{mist}{HTML}{F2F2F2}     % fond des blocs / du code

\setbeamercolor{normal text}{fg=ink, bg=white}
\setbeamercolor{alerted text}{fg=accent}
\setbeamercolor{frametitle}{bg=ink, fg=white}
\setbeamercolor{progress bar}{fg=accent, bg=mist}
\setbeamercolor{block title}{bg=mist, fg=ink}
\setbeamercolor{block body}{bg=mist, fg=ink}
\setbeamercolor{title separator}{fg=accent}

% ============================================================
%  LANGUE & ENCODAGE
% ============================================================
\usepackage[utf8]{inputenc}
\usepackage[T1]{fontenc}
\usepackage[french]{babel}
\usepackage{csquotes}
\usepackage{etoolbox}

% ============================================================
%  NOTES DE L'ORATEUR  —  deux sorties depuis une seule source
% ------------------------------------------------------------
%  Par defaut : AUCUNE note (PDF propre, celui qu'on partage).
%  Pour obtenir la version annotee, voir build.ps1 / le README.
% ============================================================
\providecommand{\NOTESMODE}{off}
\ifdefstring{\NOTESMODE}{cote}{\setbeameroption{show notes on second screen=right}}{}
\ifdefstring{\NOTESMODE}{pages}{\setbeameroption{show notes}}{}
\ifdefstring{\NOTESMODE}{seules}{\setbeameroption{show only notes}}{}
\ifdefstring{\NOTESMODE}{off}{\setbeameroption{hide notes}}{}

% mise en forme des pages de notes : gabarit minimal.
% On supprime la vignette et l'arborescence de navigation — en mode
% « cote », la diapo est deja affichee a gauche : tout l'espace va au texte.
\setbeamercolor{note page}{bg=white, fg=ink}
\setbeamertemplate{note page}{%
  \setlength{\parskip}{0.30em}%
  \setlength{\parindent}{0pt}%
  \footnotesize
  \insertnote%
  \vfill
}
% Contrainte a respecter : une note = UNE page. A cette taille, cela
% represente environ 230-240 mots, soit ~1 min 30 de parole.

% ============================================================
%  MATHS
% ============================================================
\usepackage{amsmath, amssymb, mathtools, bm}

% ============================================================
%  FIGURES & ANNOTATION DE FORMULES
% ============================================================
\usepackage{graphicx}
\usepackage{booktabs}
\usepackage{array} % pour les colonnes de tableau centrees verticalement (m{})
\usepackage{tikz}
\usetikzlibrary{tikzmark, arrows.meta, positioning, calc, fit, backgrounds}
\graphicspath{{figures/}}

% Annoter un terme d'une formule : \annot[<decalage>]{<tikzmark>}{<texte>}
\newcommand{\annot}[3][0.8cm]{%
  \begin{tikzpicture}[overlay, remember picture]
    \draw[accent, -{Stealth[length=4pt]}, thin]
      ([yshift=-#1]pic cs:#2) -- ([yshift=-2pt]pic cs:#2);
    \node[accent, font=\scriptsize, anchor=north]
      at ([yshift=-#1]pic cs:#2) {#3};
  \end{tikzpicture}%
}

% ============================================================
%  CODE
% ============================================================
\usepackage{minted}
\setminted{
  fontsize=\scriptsize,
  bgcolor=mist,
  linenos=false,
  frame=none,
  breaklines=true,
  autogobble=true,
}
\usemintedstyle{friendly}

% ============================================================
%  RACCOURCIS LTN
% ============================================================
\newcommand{\G}{\mathcal{G}}                      % grounding function
\newcommand{\dom}{\mathcal{D}}                    % domaine
\newcommand{\KB}{\mathcal{K}}                     % base de connaissances
\newcommand{\SatAgg}{\mathrm{SatAgg}}
\newcommand{\Forall}{\forall}
\newcommand{\Exists}{\exists}

% ============================================================
%  PAGE DE TITRE
% ============================================================
\title{Logic Tensor Networks}
\subtitle{Apprendre sous contrainte logique : une introduction pratique à l'IA neuro-symbolique}
\author{Léonel Junior Vodounou}
\institute{\textcolor{ash}{[institution à compléter]}}
\date{\today}

\begin{document}

\maketitle

% ============================================================
\section{Introduction / Motivation}

% ------------------------------------------------------------------
%  SLIDE 1 — Deux paradigmes
%  Idée unique : les deux paradigmes ne sont pas deux boîtes à outils
%  concurrentes, ce sont deux réponses opposées à UNE question —
%  d'où vient la structure du problème. Cet axe est celui que le
%  Bitter Lesson (slide 2) prétend trancher.
% ------------------------------------------------------------------
\begin{frame}{IA symbolique et apprentissage profond}
  \begin{center}
  \begin{tikzpicture}[x=1cm, y=1cm]

    % --- les deux extrémités de l'axe : la réponse de chaque camp -----
    \node[ash, font=\small\itshape, anchor=west] at (0, 1.55)
      {écrite par l'humain};
    \node[ash, font=\small\itshape, anchor=east] at (12, 1.55)
      {apprise depuis les données};

    % --- axe ---------------------------------------------------------
    \draw[ash, {Stealth[length=6pt]}-{Stealth[length=6pt]}, semithick]
      (0, 0.65) -- (12, 0.65);

    % --- pôle gauche --------------------------------------------------
    \fill[accent] (2.3, 0.65) circle (2.6pt);
    \node[ink, font=\bfseries, anchor=south] at (2.3, 0.88)
      {IA symbolique};
    \node[ink, font=\footnotesize, anchor=north, text width=4.6cm, align=center]
      at (2.3, 0.32) {\textbf{+} transparente, vérifiable,\\ très économe en données};
    \node[ash, font=\footnotesize, anchor=north, text width=4.6cm, align=center]
      at (2.3, -0.75) {\textbf{--} rigide, fragile au bruit,\\ ne passe pas à l'échelle};

    % --- pôle droit ---------------------------------------------------
    \fill[accent] (9.7, 0.65) circle (2.6pt);
    \node[ink, font=\bfseries, anchor=south] at (9.7, 0.88)
      {Apprentissage profond};
    \node[ink, font=\footnotesize, anchor=north, text width=4.6cm, align=center]
      at (9.7, 0.32) {\textbf{+} passe à l'échelle,\\ apprend depuis le signal brut};
    \node[ash, font=\footnotesize, anchor=north, text width=4.6cm, align=center]
      at (9.7, -0.75) {\textbf{--} opaque, invérifiable,\\ très gourmand en données};

  \end{tikzpicture}
  \end{center}

  \vspace{0.3em}
  \begin{center}
    \small Les deux limites sont \alert{symétriques} : chacun paie sa force
    par l'exacte faiblesse de l'autre.
  \end{center}
\end{frame}

\note{%
\textbf{\textcolor{accent}{Slide 1 : IA symbolique et apprentissage profond}} \hfill \textit{$\approx$ 1 min 15}

Avant d'examiner les capacités d'une intelligence artificielle, une question
plus fondamentale se pose : \textit{sur quoi repose sa connaissance du monde,
et comment celle-ci est-elle construite ?}

Et à cette question, deux grandes approches existent.

Soit on lui transmet explicitement des connaissances : autrement on écrit les
règles, on encode la logique. C'est l'IA symbolique.

Soit on lui fournit des données et on la laisse apprendre des régularités à
partir de celles-ci : c'est l'apprentissage profond.

Mais chacune de ces approches a ses limites.

Le symbolique est transparent : on peut suivre les règles qui conduisent à une
conclusion. En revanche, il reste rigide : dès qu'une situation s'écarte de ce
qui a été prévu, ses règles peuvent devenir insuffisantes.

Le réseau de neurones est beaucoup plus flexible : il apprend directement à
partir des données et peut s'adapter à des situations complexes. Mais cette
capacité a un coût : son fonctionnement est beaucoup plus difficile à
interpréter.

Autrement dit, les deux approches échouent sur des points presque opposés.

Alors une question se pose : \textit{peut-on réunir leurs forces sans hériter
de leurs faiblesses ?}

C'est précisément le problème auquel nous allons nous intéresser.

\smallskip
\textcolor{ash}{\textit{Si la question est posée :} cette opposition structure le
champ depuis les années 1980.}
}

% ------------------------------------------------------------------
%  SLIDE 2 — Le Bitter Lesson
%  Idée unique : la connaissance codée à la main plafonne ; le calcul,
%  lui, continue de croître (loi de Moore) et finit par la dépasser.
%  Le graphe porte cette affirmation directement, sans texte.
% ------------------------------------------------------------------
\begin{frame}{La Bitter Lesson}
  \begin{center}
  \begin{tikzpicture}[x=0.92cm, y=0.72cm]

    % --- axes ----------------------------------------------------
    \draw[ash, -{Stealth[length=6pt]}, semithick]
      (0, 0) -- (11.3, 0) node[right, font=\footnotesize, ash] {calcul disponible};
    \draw[ash, -{Stealth[length=6pt]}, semithick]
      (0, 0) -- (0, 6.3) node[above, font=\footnotesize, ash] {performance};

    % --- courbe : connaissance codee a la main (plafonne) ---------
    \draw[ash, thick]
      plot [smooth] coordinates {(0.3,0.4) (1,1.6) (2,2.6) (3,3.1)
                                  (4,3.3) (6,3.35) (8,3.3) (10.3,3.3)};
    \node[ash, font=\footnotesize\itshape, anchor=west]
      at (8.5, 3.55) {connaissance codée à la main};

    % --- courbe : calcul, recherche et apprentissage (continue) ---
    \draw[accent, thick]
      plot [smooth] coordinates {(0.3,0.15) (1,0.35) (2,0.7) (3,1.2)
                                  (4,1.9) (5,2.7) (6,3.4) (7,4.0)
                                  (8,4.6) (9,5.1) (10.3,5.6)};
    \node[accent, font=\footnotesize\bfseries, anchor=south west]
      at (6.3, 5.75) {calcul : recherche et apprentissage};

  \end{tikzpicture}
  \end{center}

  \vspace{0.1em}
  \begin{center}
    \small Ce que le calcul apprend seul finit par \alert{dépasser} ce que
    nous codons à la main.
  \end{center}
\end{frame}

\note{%
\textbf{\textcolor{accent}{Slide 2 : La Bitter Lesson}} \hfill \textit{$\approx$ 1 min 30}

En 2019, le chercheur Richard Sutton publie un texte devenu une référence
dans le domaine, intitulé la Bitter Lesson.

Sa thèse tient en une phrase. Sur soixante dix ans de recherche en
intelligence artificielle, les méthodes générales qui exploitent le calcul
finissent toujours par surpasser celles qui exploitent la connaissance
humaine, et l'écart est large.

La raison tient à la loi de Moore. Le calcul disponible croît sans cesse,
alors qu'une connaissance codée à la main reste fixe. À court terme,
injecter de la connaissance humaine aide, et cela satisfait le chercheur.
Mais à long terme, cette connaissance finit par limiter le système : elle ne
sait pas exploiter le calcul supplémentaire qui devient disponible.

L'exemple le plus connu reste les échecs. En 1997, Deep Blue bat Kasparov
grâce à une recherche massive, et non grâce à une compréhension fine du jeu.
Le même motif s'est ensuite répété au jeu de Go, en reconnaissance vocale,
en vision par ordinateur. À chaque fois, une méthode générale, appuyée sur
le calcul, finit par dépasser des années de connaissance construite à la
main.

Sutton en tire deux leçons. Il faut privilégier les méthodes qui profitent du
calcul supplémentaire à mesure qu'il devient disponible. Et il faut cesser de
vouloir coder directement ce que nous croyons savoir du monde, car cette
complexité est sans fin.

\smallskip
\textcolor{ash}{\textit{Si la question est posée :} le texte complet tient en
deux pages et circule très largement dans la communauté depuis 2019.}
}

% ------------------------------------------------------------------
%  SLIDE 3 — Le compromis neuro-symbolique
%  Idée unique : LTN ne fixe pas la réponse, il fixe le vocabulaire
%  dans lequel le gradient va la chercher. La connaissance humaine
%  porte sur la grammaire, pas sur les valeurs. C'est la réponse
%  directe à l'objection de Sutton, avec sa limite assumée.
% ------------------------------------------------------------------
\begin{frame}{Le compromis neuro-symbolique}
  \begin{center}
  \begin{tikzpicture}[x=1cm, y=1cm]

    % --- boites : dessinees separement du texte -------------------
    % (un noeud qui combine a la fois "draw" et "text width" force pgf
    %  a mesurer le texte en deux passes ; en mode notes, cette shape
    %  laisse parfois une trace de son contour sur la page suivante.
    %  On separe donc le rectangle, simple trace, du texte, simple bloc.)
    \draw[ash, thick, rounded corners] (0, -1.7) rectangle (5.4, 1.7);
    \draw[accent, thick, rounded corners] (6.8, -1.7) rectangle (12.2, 1.7);

    % --- contenu de la boite gauche : ce que nous fixons -----------
    \node[align=center, text width=4.7cm] at (2.7, 0) {
        \textcolor{ash}{\scshape\footnotesize fixé par nous}\\[5pt]
        \textbf{le vocabulaire logique}\\[6pt]
        \textcolor{ash}{\footnotesize prédicats, connecteurs,\\ quantificateurs}\\[5pt]
        $P(x),\ \land,\ \lor,\ \forall,\ \exists$
      };

    % --- contenu de la boite droite : ce que le gradient trouve ----
    \node[align=center, text width=4.7cm] at (9.5, 0) {
        \textcolor{accent}{\scshape\footnotesize trouvé par le gradient}\\[5pt]
        \textbf{les paramètres du réseau}\\[6pt]
        \textcolor{ash}{\footnotesize appris par descente de gradient,\\
        depuis les données}\\[5pt]
        $\theta$
      };

    % --- fleche -----------------------------------------------------
    \draw[ash, {Stealth[length=7pt]}-, thick] (6.7, 0) -- (5.5, 0)
      node[midway, above, font=\footnotesize\itshape, ash] {contraint};

  \end{tikzpicture}
  \end{center}

  \vspace{0.4em}
  \begin{center}
    \small Les LTN fixent la forme d'un prédicat, mais pas \alert{la
    fonction qu'il apprend à représenter}.
  \end{center}
\end{frame}

\note{%
\textbf{\textcolor{accent}{Slide 3 : Le compromis neuro-symbolique}} \hfill \textit{$\approx$ 1 min 45}

Le neuro symbolique semble à première vue contredire directement cette leçon.
Il propose d'injecter de la logique, une connaissance humaine explicite, dans
un système d'apprentissage.

En effet, ici, il faut distinguer deux niveaux.

Le premier est celui de la perception : la fonction qui transforme une
donnée brute en une estimation de ce qu'elle représente. Rien n'y est codé à
l'avance, elle reste un réseau de neurones ordinaire, appris entièrement par
descente de gradient. Sur ce plan, LTN suit la leçon de Sutton : une méthode
générale qui exploite le calcul pour apprendre.

Le second niveau est celui du raisonnement : la façon dont plusieurs
perceptions s'articulent pour rendre compte d'une observation globale. Là,
nous intervenons directement, en écrivant une formule logique, avec des
\textbf{prédicats, des connecteurs et des quantificateurs}. Je reviendrai sur
chacun de ces termes juste après. Cette structure n'est pas apprise, elle est
injectée, comme un biais explicite.

LTN occupe donc une position intermédiaire. Il applique le principe de
Sutton à l'étage perceptif, mais s'en écarte à l'étage symbolique, où une
connaissance humaine reste programmée à la main. C'est un compromis assumé :
nous renonçons à une part de généralité en échange d'une meilleure
efficacité statistique, surtout utile quand les données sont rares.

C'est ce fonctionnement des Logic Tensor Networks que nous allons maintenant
explorer dans la suite de cette présentation.

\smallskip
\textcolor{ash}{\textit{Si la question est posée :} cette défense a une
limite reconnue, elle sera reprise en conclusion.}
}

% ============================================================
\section{Fondements théoriques}

% ------------------------------------------------------------------
%  SLIDE 4 — Le vocabulaire de la logique classique
%  Idee unique : quatre briques suffisent a ecrire une formule
%  (domaine, constante, variable, predicat). Tient la promesse
%  faite en slide 3 ("je reviendrai sur chacun de ces termes").
%  Les connecteurs et quantificateurs viennent a la diapo suivante :
%  ce sont eux qui COMBINENT ces briques, pas des briques elles-memes.
% ------------------------------------------------------------------
\begin{frame}{Logique du premier ordre : le vocabulaire}
  \begin{center}
    \small\itshape\color{ash}
    « Socrate est un homme. » \quad « Tous les hommes sont mortels. »\\
    \normalfont\small\normalcolor
    un langage formel : des phrases précises, sans ambiguïté, vérifiables.
  \end{center}

  \vspace{0.7em}
  \begin{center}
  \small
  \renewcommand{\arraystretch}{1.6}
  \begin{tabular}{@{} l c >{\raggedright\arraybackslash}m{5.6cm} @{}}
    \toprule
    {\color{ash}\scshape\footnotesize notion} &
    {\color{ash}\scshape\footnotesize notation} &
    {\color{ash}\scshape\footnotesize exemple} \\
    \midrule
    domaine & $\mathcal{D}$ & $\{$Léonel, Mélène, Alice, Abdel$\}$ \\
    constante & $\text{abdel} \in \mathcal{D}$ & pointe toujours vers Abdel \\
    variable & $x$ & un individu quelconque du domaine \\
    prédicat & $\text{EstÉtudiant}(x)$ & vrai pour Abdel, faux pour Mélène \\
    \bottomrule
  \end{tabular}
  \end{center}

  \vspace{0.5em}
  \begin{center}
    \small Ces quatre briques suffisent à écrire n'importe quelle
    \alert{formule}.
  \end{center}
\end{frame}

\note{%
\textbf{\textcolor{accent}{Slide 4 : Le vocabulaire de la logique classique}} \hfill \textit{$\approx$ 1 min}

Avant de voir comment Logic Tensor Networks fonctionne, reprenons ce qu'est
la logique du premier ordre elle même.

C'est un langage formel : un ensemble de règles d'écriture strictes, qui
permet de décrire le monde par des phrases précises, sans ambiguïté et
vérifiables. « Socrate est un homme » ou « tous les hommes sont mortels »
sont des phrases de ce type.

Ce langage repose sur quatre briques. Le domaine est l'ensemble des objets
dont on parle. Si je construis un système pour parler d'un petit groupe de
personnes, mon domaine pourrait être Léonel, Mélène, Alice et Abdel.

Une constante désigne toujours le même objet du domaine : le symbole abdel
pointe vers Abdel, et rien d'autre.

Une variable, notée x, désigne un objet quelconque du domaine, non encore
fixé.

Un prédicat exprime une propriété, vraie ou fausse selon l'objet auquel on
l'applique. Par exemple, EstÉtudiant appliqué à Abdel est vrai, alors
qu'appliqué à Mélène il est faux.

Ces quatre notions suffisent à écrire n'importe quelle formule. Mais un
prédicat peut prendre plus d'un argument, et il existe aussi les fonctions.
C'est l'objet de la diapositive suivante.
}

% ------------------------------------------------------------------
%  SLIDE 5 — Predicats (arite) et fonctions
%  Idee unique : un predicat peut prendre plusieurs arguments (arite),
%  et une fonction se distingue d'un predicat par ce qu'elle renvoie
%  (un objet du domaine, pas une verite). Reprend fidelement les
%  sections 5, 6 et 7 de la cellule "Prerequis : logique du premier
%  ordre" (unaire avant binaire, arite, D^n, puis fonctions).
% ------------------------------------------------------------------
\begin{frame}{Prédicats et fonctions}
  \begin{center}
    \small\itshape\color{ash}
    l'arité d'un prédicat est son nombre d'arguments : $n$ arguments
    $\Rightarrow$ prédicat $n$-aire, défini sur $\mathcal{D}^n$.
  \end{center}

  \vspace{0.6em}
  \begin{center}
  \footnotesize
  \renewcommand{\arraystretch}{1.7}
  \begin{tabular}{@{} l c >{\raggedright\arraybackslash}m{3.0cm} @{}}
    \toprule
    {\color{ash}\scshape\footnotesize notion} &
    {\color{ash}\scshape\footnotesize notation} &
    {\color{ash}\scshape\footnotesize exemple} \\
    \midrule
    prédicat unaire (1-aire) & $\text{EstÉtudiant}(x)$ &
      vrai pour Abdel, faux pour Mélène \\
    prédicat binaire (2-aire) & $\text{EstAmoureuxDe}(x,y)$ &
      vrai pour (Alice, Abdel) \\
    fonction & $\text{PèreDe}(x)$ &
      $\text{PèreDe}(\text{Alice}) = \text{Léonel}$ \\
    \bottomrule
  \end{tabular}
  \end{center}

  \vspace{0.6em}
  \begin{center}
    \small Un prédicat renvoie une \alert{vérité} ; une fonction renvoie un
    objet du domaine.
  \end{center}
\end{frame}

\note{%
\textbf{\textcolor{accent}{Slide 5 : Prédicats et fonctions}} \hfill \textit{$\approx$ 2 min}

Reprenons les prédicats, avec toute leur richesse.

Un prédicat est une propriété ou une relation qu'on peut attribuer à un ou
plusieurs objets, et qui est vraie ou fausse.

Prenons d'abord un prédicat à un seul argument, dit unaire, ou un-aire.
EstÉtudiant de x prend un seul objet en entrée et dit si cet objet a la
propriété d'être étudiant. EstÉtudiant appliqué à Abdel vaut vrai ; appliqué
à Mélène, il vaut faux.

Un prédicat peut aussi prendre plusieurs arguments. EstAmoureuxDe de x et y,
un prédicat binaire, ou deux-aire, exprime une relation entre deux objets.
EstAmoureuxDe appliqué à Alice et Abdel vaut vrai. Rien ne garantit que la
relation soit symétrique : EstAmoureuxDe appliqué à Abdel et Alice pourrait
très bien avoir une valeur différente.

Le nombre d'arguments d'un prédicat s'appelle son arité. Un prédicat à $n$
arguments est dit $n$-aire, et il est défini sur $\mathcal{D}^n$,
c'est-à-dire l'ensemble de tous les $n$-uplets d'objets du domaine.

Les fonctions, elles, sont différentes. Contrairement à un prédicat, une
fonction ne renvoie pas une valeur de vérité, mais un autre objet du
domaine. PèreDe de x, appliquée à une personne, renvoie une autre personne,
son père. Avec notre domaine, PèreDe appliqué à Alice vaut Léonel : la
sortie reste dans le domaine.

Nous avons maintenant tout le vocabulaire nécessaire : domaine, constantes,
variables, prédicats de toute arité, et fonctions. Voyons à présent comment
les combiner.
}

% ------------------------------------------------------------------
%  SLIDE 6 — Connecteurs et quantificateurs
%  Idee unique : ces particules combinent les briques de la slide 4
%  en formules completes. Le paiement de la diapo est la formule
%  finale elle-meme (meme domaine : Leonel, Melene, Alice, Abdel).
% ------------------------------------------------------------------
\begin{frame}{Connecteurs et quantificateurs}
  \begin{center}
  \small
  \renewcommand{\arraystretch}{1.5}
  % Une seule tabulation (pas de columns/minipage cote a cote : ce
  % type de construction a montre une fuite de bordure sur la page
  % de notes adjacente sur des diapos precedentes).
  \begin{tabular}{@{} c l @{\hspace{1.4cm}} c l @{}}
    \toprule
    {\color{ash}\scshape\footnotesize symbole} &
    {\color{ash}\scshape\footnotesize connecteur} &
    {\color{ash}\scshape\footnotesize symbole} &
    {\color{ash}\scshape\footnotesize quantificateur} \\
    \midrule
    $\lnot$ & négation (non) & $\forall x\,\phi(x)$ & pour tout $x$ \\
    $\land$ & conjonction (et) & $\exists x\,\phi(x)$ & il existe un $x$ \\
    $\lor$ & disjonction (ou) & & \\
    $\implies$ & implication (si\ldots alors) & & \\
    \bottomrule
  \end{tabular}
  \end{center}

  \vspace{1.1em}
  \begin{center}
    \small Ensemble, ils permettent d'écrire une \alert{formule} complète :
  \end{center}
  \[
    \forall x\,\bigl(\text{EstÉtudiant}(x) \implies \exists y\,
    \text{EstAmiDe}(x,y)\bigr)
  \]
  \begin{center}
    \small\itshape\color{ash} « tout étudiant a au moins un ami »
  \end{center}
\end{frame}

\note{%
\textbf{\textcolor{accent}{Slide 6 : Connecteurs et quantificateurs}} \hfill \textit{$\approx$ 1 min 15}

Les connecteurs permettent de combiner plusieurs propositions en une seule.
La négation ($\lnot$) inverse une valeur de vérité. La conjonction ($\land$)
n'est vraie que si les deux propositions le sont ensemble. La disjonction
($\lor$) est vraie dès que l'une des deux l'est. L'implication ($\implies$)
exprime que si la première proposition est vraie, la seconde l'est
nécessairement.

Les quantificateurs, quant à eux, permettent de parler de plusieurs objets
sans les nommer un par un. Pour tout x, phi de x ($\forall x\,\phi(x)$),
affirme que la propriété vaut pour tout le domaine. Il existe un x, phi de x
($\exists x\,\phi(x)$), affirme qu'elle vaut pour au moins un objet.

Une fois tout cela en main, on peut écrire une formule complète. Pour tout
x, si x est étudiant, alors il existe un y tel que x est ami avec y
($\forall x\,(\text{EstÉtudiant}(x) \implies \exists y\,
\text{EstAmiDe}(x,y))$). Autrement dit, tout étudiant a au moins un ami.

Cette formule reste, pour l'instant, entièrement discrète : vraie ou fausse,
sans nuance possible. C'est précisément ce que nous allons chercher à
dépasser dans la suite.
}

% ------------------------------------------------------------------
%  SLIDE 7 — L'interpretation classique I
%  Idee unique : I est un choix arbitraire parmi d'autres. La meme
%  formule, sous une autre interpretation, dit autre chose. Ce
%  constat motive directement le terme suivant : le grounding.
% ------------------------------------------------------------------
\begin{frame}{L'interprétation : le pont entre syntaxe et sens}
  \begin{center}
  \footnotesize
  \renewcommand{\arraystretch}{1.7}
  \begin{tabular}{@{} l c >{\raggedright\arraybackslash}m{3.0cm} @{}}
    \toprule
    {\color{ash}\scshape\footnotesize symbole} &
    {\color{ash}\scshape\footnotesize interprétation $\mathcal{I}$} &
    {\color{ash}\scshape\footnotesize exemple} \\
    \midrule
    constante $c$ & $\mathcal{I}(c) \in \mathcal{D}$ &
      $\mathcal{I}(\text{abdel}) = \text{Abdel}$ \\
    prédicat $P$ ($n$-aire) & $\mathcal{I}(P) \subseteq \mathcal{D}^n$ &
      $\mathcal{I}(\text{EstÉtudiant})$: Alice, Abdel \\
    fonction $f$ & $\mathcal{I}(f) : \mathcal{D}^n \to \mathcal{D}$ &
      $\mathcal{I}(\text{PèreDe})(\text{Alice}) = \text{Léonel}$ \\
    \bottomrule
  \end{tabular}
  \end{center}

  \vspace{0.7em}
  \begin{center}
    \small Rien n'oblige à choisir \alert{cette} interprétation plutôt
    qu'une autre.
  \end{center}
\end{frame}

\note{%
\textbf{\textcolor{accent}{Slide 7 : L'interprétation}} \hfill \textit{$\approx$ 1 min 30}

L'interprétation $\mathcal{I}$ fait le pont entre la syntaxe et la
sémantique. Elle assigne à chaque symbole brut du langage un objet concret,
dans un domaine choisi au préalable.

Ce n'est pas une règle unique, mais un ensemble de trois assignations. À
chaque constante, elle associe un élément précis du domaine. À chaque
prédicat $n$-aire, elle associe un sous-ensemble de $\mathcal{D}^n$ :
exactement les combinaisons pour lesquelles il est vrai. À chaque fonction,
elle associe une véritable fonction mathématique sur le domaine.

Prenons EstÉtudiant. Si $\mathcal{I}(\text{EstÉtudiant})$ vaut Alice et
Abdel, alors évaluer EstÉtudiant appliqué à Abdel revient à demander si
Abdel appartient à cet ensemble. C'est le cas, donc la formule est vraie.

Mais voici le point crucial. Rien n'oblige à choisir cette interprétation
plutôt qu'une autre. Reprenons exactement les mêmes symboles, abdel et
EstÉtudiant, mais changeons complètement d'interprétation. Choisissons cette
fois un domaine de nombres, et décidons que abdel désigne le nombre 5, et
qu'EstÉtudiant désigne la propriété d'être pair. La même formule
syntaxique, EstÉtudiant appliqué à abdel, devient alors une affirmation
totalement différente : est-ce que 5 est pair ? Ce qui est évidemment faux.
Ainsi, la syntaxe n'a pas changé d'un seul symbole, mais son sens s'est
complètement inversé.

C'est précisément ce point qui va nous conduire au terme suivant : le
grounding.
}

% ------------------------------------------------------------------
%  SLIDE 8 — Real Logic : le grounding G
%  Idee unique : G n'est pas une interpretation quelconque, elle est
%  contrainte a cibler l'espace des tenseurs. Deux cibles seulement :
%  terme -> tenseur, formule -> reel dans [0,1]. C'est ce remplacement
%  qui rend la logique differentiable (cf. slide 3, meme mecanisme).
% ------------------------------------------------------------------
\begin{frame}{Real Logic : le grounding $\mathcal{G}$}
  \begin{center}
  \footnotesize
  \renewcommand{\arraystretch}{1.8}
  \begin{tabular}{@{} l >{\centering\arraybackslash}m{2.8cm}
                       >{\centering\arraybackslash}m{2.8cm} @{}}
    \toprule
    & {\color{ash}\scshape\footnotesize interprétation $\mathcal{I}$} &
      {\color{accent}\scshape\footnotesize grounding $\mathcal{G}$} \\
    \midrule
    terme & un objet abstrait de $\mathcal{D}$ & un tenseur de réels \\
    formule & vrai ou faux ($\{0,1\}$) & un réel dans $[0,1]$ \\
    \bottomrule
  \end{tabular}
  \end{center}

  \vspace{0.7em}
  \begin{center}
    \small $\mathcal{G}$ n'est pas une interprétation quelconque : elle est
    \alert{contrainte} à l'espace des tenseurs.
  \end{center}
\end{frame}

\note{%
\textbf{\textcolor{accent}{Slide 8 : Real Logic}} \hfill \textit{$\approx$ 1 min 45}

Nous avons vu que rien ne contraint une interprétation classique. Real
Logic, la sémantique utilisée par Logic Tensor Networks, résout ce problème
d'une manière particulière. Elle ne choisit pas une interprétation parmi
d'autres : elle en contraint la forme.

Le domaine n'est plus un ensemble abstrait. Il est interprété concrètement
par des tenseurs réels. C'est pour cette raison qu'on parle de grounding,
noté $\mathcal{G}$, plutôt que d'interprétation : ce nouveau terme marque
que $\mathcal{G}$ n'est pas une interprétation quelconque, mais une
interprétation contrainte à toujours cibler l'espace des tenseurs.

Concrètement, $\mathcal{G}$ agit sur deux catégories de symboles. À tout
terme, une constante, une variable, le résultat d'une fonction, $\mathcal{G}$
associe un tenseur de réels. À toute formule, $\mathcal{G}$ associe un réel
dans l'intervalle $[0,1]$ : un degré de vérité continu, et non plus un
booléen strict.

Reprenons EstÉtudiant appliqué à abdel. Sous $\mathcal{I}$, la réponse était
vraie ou fausse. Sous $\mathcal{G}$, la même expression devient un nombre,
disons $0{,}92$ : proche de la vérité, sans y être totalement.

C'est ce remplacement, des objets abstraits par des tenseurs, du booléen par
le continu, qui rend une connaissance logique exploitable comme contrainte
différentiable, entraînable par descente de gradient.
}

% ------------------------------------------------------------------
%  SLIDE 9 — Grounding concret : constantes, predicats, fonctions,
%  variables
%  Idee unique : G n'est pas qu'une definition abstraite, elle se
%  code tres concretement. Un predicat = reseau + sigmoide ; une
%  fonction = reseau sans sigmoide ; c'est le seul signal qui les
%  distingue. Une constante peut etre fixe ou apprenable (rappelle
%  la slide 3). Une variable est un batch, pas un individu.
% ------------------------------------------------------------------
\begin{frame}{Grounder les symboles non logiques}
  \begin{center}
  \footnotesize
  \renewcommand{\arraystretch}{1.55}
  \begin{tabular}{@{} l >{\raggedright\arraybackslash}m{6.4cm} @{}}
    \toprule
    {\color{ash}\scshape\footnotesize symbole} &
    {\color{ash}\scshape\footnotesize grounding} \\
    \midrule
    constante & tenseur fixé (donnée) ou \alert{apprenable} (embedding,
      appris par gradient) \\
    prédicat & fonction (souvent un réseau) se terminant par une
      \textbf{sigmoïde} : sortie dans $[0,1]$ \\
    fonction & fonction (souvent un réseau) \textbf{sans} sigmoïde :
      sortie un tenseur quelconque \\
    variable & un \textbf{batch}, ou séquence, d'individus \\
    \bottomrule
  \end{tabular}
  \end{center}

  \vspace{0.3em}
  \begin{center}
    \small La \alert{sigmoïde} est le seul signal qui distingue un
    prédicat d'une fonction.
  \end{center}
\end{frame}

\note{%
\textbf{\textcolor{accent}{Slide 9 : Grounder les symboles non logiques}} \hfill \textit{$\approx$ 1 min 15}

Passons maintenant du principe général à sa mise en œuvre concrète.

Une constante est groundée sur un tenseur réel : fixé à partir d'une
donnée, ou entraînable comme un embedding, ajusté par descente de gradient.
Nous retrouvons ici, très concrètement, la distinction posée dès la slide
sur le compromis neuro-symbolique.

Un prédicat associe à son entrée un réel dans $[0,1]$ : une fonction
mathématique quelconque, souvent un réseau de neurones, dont la dernière
couche se termine par une sigmoïde. Une fonction, elle, n'a aucune
contrainte de sortie : jamais de sigmoïde. C'est ce seul détail qui
distingue les deux, dans n'importe quel notebook LTN.

Une variable, enfin, représente un batch, ou séquence, d'individus. C'est
cette notion qui permettra d'écrire des quantifications comme pour tout
$x$, puisque cela suppose que $x$ parcoure plusieurs individus.
}

% ------------------------------------------------------------------
%  SLIDE 10 — Connecteurs flous : le choix du produit
%  Idee unique : les connecteurs deviennent des fonctions lisses sur
%  [0,1], et le choix precis de la famille (produit plutot que
%  Godel min/max) n'est pas arbitraire : il conditionne si le
%  gradient circule sur les deux arguments ou un seul.
% ------------------------------------------------------------------
\begin{frame}{La configuration Product Real Logic}
  \begin{center}
  \footnotesize
  \renewcommand{\arraystretch}{1.7}
  \begin{tabular}{@{} >{\centering\arraybackslash}p{2.0cm}
                       >{\centering\arraybackslash}p{2.2cm}
                       >{\raggedright\arraybackslash}p{3.4cm} @{}}
    \toprule
    {\color{ash}\scshape\footnotesize connecteur} &
    {\color{ash}\scshape\footnotesize opérateur} &
    {\color{ash}\scshape\footnotesize nom} \\
    \midrule
    $\lnot u$ & $1-u$ & négation standard \\
    $u \land v$ & $uv$ & t-norme produit \\
    $u \lor v$ & $u+v-uv$ & t-conorme produit \\
    $u \implies v$ & $1-u+uv$ & implication de Reichenbach \\
    \bottomrule
  \end{tabular}
  \end{center}

  \vspace{0.7em}
  \begin{center}
    \small Le produit répartit le \alert{gradient} sur les deux arguments ;
    le min/max de Gödel, non.
  \end{center}
\end{frame}

\note{%
\textbf{\textcolor{accent}{Slide 10 : La configuration Product Real Logic}} \hfill \textit{$\approx$ 1 min 30}

Les connecteurs logiques classiques ne sont définis que pour des valeurs
strictement zéro ou un. Logic Tensor Networks les remplace donc par des
opérateurs de logique floue, définis sur tout l'intervalle continu $[0,1]$.

Plusieurs familles d'opérateurs flous existent, mais elles ne sont pas
équivalentes du point de vue de l'entraînement. Le minimum et le maximum de
Gödel ($u \land v = \min(u,v)$, $u \lor v = \max(u,v)$), par exemple, ne
font circuler le gradient que sur un seul argument à la fois, ce qui peut
ralentir ou bloquer l'apprentissage.

C'est pourquoi Logic Tensor Networks recommande en général la
configuration Product Real Logic. Pour la négation, la négation standard :
$1-u$. Pour la conjonction, la t-norme produit : $u \land v = uv$. Pour la
disjonction, la t-conorme produit : $u \lor v = u+v-uv$. Et pour
l'implication, l'implication de Reichenbach : $u \implies v = 1-u+uv$.

Ce choix n'est pas arbitraire. Contrairement au minimum et au maximum, le
produit a un gradient qui se répartit sur les deux arguments à la fois, ce
qui en fait un choix numériquement plus stable pour l'entraînement.

Chaque connecteur devient ainsi une fonction lisse sur $[0,1]$, dérivable
presque partout. C'est cette propriété qui permet à la descente de gradient
de fonctionner à travers toute la formule.

\smallskip
\textcolor{ash}{\textit{Si la question est posée :} il existe d'autres
familles, Gödel ou Łukasiewicz, mais elles posent des problèmes de gradient
détaillés dans un notebook complémentaire.}
}

% ------------------------------------------------------------------
%  SLIDE 11 — Quantificateurs : la moyenne generalisee
%  Idee unique : pMean et pMeanError sont des versions lissees du
%  max et du min, pilotees par p. p grand = fidelite logique stricte,
%  p proche de 1 = tolerance au bruit, meilleur gradient.
% ------------------------------------------------------------------
\begin{frame}{Quantificateurs : la moyenne généralisée}
  \[
    \exists x\,\phi(x) \;\approx\; \mathrm{pMean} =
    \Big(\tfrac{1}{n}\textstyle\sum_{i=1}^n u_i^p\Big)^{1/p}
  \]
  \vspace{-0.3em}
  \[
    \forall x\,\phi(x) \;\approx\; \mathrm{pMeanError} =
    1-\Big(\tfrac{1}{n}\textstyle\sum_{i=1}^n (1-u_i)^p\Big)^{1/p}
  \]

  \vspace{0.3em}
  \begin{itemize}
    \small
    \item quand $p \to 1$ : les deux tendent vers une moyenne simple,
      tolérante aux cas difficiles du batch ;
    \item quand $p \to \infty$ : $\mathrm{pMean}$ tend vers le
      \alert{maximum}, $\mathrm{pMeanError}$ vers le \alert{minimum}\\
      strict, exigeant.
  \end{itemize}
\end{frame}

\note{%
\textbf{\textcolor{accent}{Slide 11 : Quantificateurs}} \hfill \textit{$\approx$ 1 min 30}

Le quantificateur universel et le quantificateur existentiel ne peuvent pas
non plus être utilisés dans leur forme classique. Ils supposent de vérifier
une propriété sur un domaine entier, alors qu'en Logic Tensor Networks, une
variable ne représente qu'un batch fini d'individus, avec des degrés de
vérité continus.

LTN les remplace donc par des opérateurs d'agrégation, qui condensent les
degrés de vérité d'un batch en une seule valeur. Pour l'existentiel, la
moyenne généralisée des degrés de vérité, notée pMean. Pour l'universel, la
moyenne généralisée
des écarts à la vérité, notée pMeanError.

Le paramètre p règle un compromis. Quand p tend vers l'infini,
pMeanError se rapproche du vrai minimum, et le quantificateur universel
retrouve son sens logique strict : un seul individu défaillant suffit à
faire chuter tout le résultat. Symétriquement, pMean se rapproche du vrai
maximum, et l'existentiel redevient facile à satisfaire.

Quand p se rapproche de 1, les deux opérateurs se rapprochent au contraire
d'une simple moyenne. Un individu défaillant peut être compensé par
plusieurs bons individus. C'est un écart volontaire par rapport à la
sémantique stricte, utile pour absorber le bruit dans les données, ou pour
obtenir un gradient mieux réparti à l'entraînement.

\smallskip
\textcolor{ash}{\textit{Si la question est posée :} la démonstration
rigoureuse de cette convergence se trouve dans le notebook
\texttt{2-grounding\_connectives.ipynb}.}
}

% ------------------------------------------------------------------
%  SLIDE 12 — Quantification diagonale
%  Idee unique : par defaut, deux variables sont croisees (toutes les
%  combinaisons). ltn.diag force une correspondance un a un, comme un
%  zip. Indispensable des qu'on a des paires (donnee, label).
% ------------------------------------------------------------------
\begin{frame}{Quantification diagonale}
  \begin{center}
    \small Par défaut, un prédicat sur deux variables est évalué sur
    \alert{toutes les combinaisons}.
  \end{center}

  \vspace{0.1em}
  \begin{center}
  \begin{tikzpicture}[x=1.05cm, y=1.05cm]

    % --- en-tetes de colonnes (labels l) et de lignes (points x) -------
    \node[ash, font=\small] at (0, 2.75) {$l_0$};
    \node[ash, font=\small] at (1, 2.75) {$l_1$};
    \node[ash, font=\small] at (2, 2.75) {$l_2$};
    \node[ash, font=\small, anchor=east] at (-0.35, 2) {$x_0$};
    \node[ash, font=\small, anchor=east] at (-0.35, 1) {$x_1$};
    \node[ash, font=\small, anchor=east] at (-0.35, 0) {$x_2$};

    % --- ligne x_0 (y=2) : diagonale en (l_0), croix ailleurs ----------
    \fill[accent!12] (-0.45,1.55) rectangle (0.45,2.45);
    \draw[accent, thick] (-0.45,1.55) rectangle (0.45,2.45);
    \node[accent, font=\footnotesize] at (0,2) {$(x_0,l_0)$};
    \draw[ash!45, thin] (0.55,1.55) rectangle (1.45,2.45);
    \node[ash!70] at (1,2) {$\times$};
    \draw[ash!45, thin] (1.55,1.55) rectangle (2.45,2.45);
    \node[ash!70] at (2,2) {$\times$};

    % --- ligne x_1 (y=1) : diagonale en (l_1), croix ailleurs ----------
    \draw[ash!45, thin] (-0.45,0.55) rectangle (0.45,1.45);
    \node[ash!70] at (0,1) {$\times$};
    \fill[accent!12] (0.55,0.55) rectangle (1.45,1.45);
    \draw[accent, thick] (0.55,0.55) rectangle (1.45,1.45);
    \node[accent, font=\footnotesize] at (1,1) {$(x_1,l_1)$};
    \draw[ash!45, thin] (1.55,0.55) rectangle (2.45,1.45);
    \node[ash!70] at (2,1) {$\times$};

    % --- ligne x_2 (y=0) : diagonale en (l_2), croix ailleurs ----------
    \draw[ash!45, thin] (-0.45,-0.45) rectangle (0.45,0.45);
    \node[ash!70] at (0,0) {$\times$};
    \draw[ash!45, thin] (0.55,-0.45) rectangle (1.45,0.45);
    \node[ash!70] at (1,0) {$\times$};
    \fill[accent!12] (1.55,-0.45) rectangle (2.45,0.45);
    \draw[accent, thick] (1.55,-0.45) rectangle (2.45,0.45);
    \node[accent, font=\footnotesize] at (2,0) {$(x_2,l_2)$};

  \end{tikzpicture}
  \end{center}

  \vspace{0.3em}
  \begin{center}
    \small Mais avec la \alert{quantification diagonale}, seules les
    paires en correspondance un à un sont conservées : les $\times$
    n'ont aucun sens pour un dataset supervisé.
  \end{center}
\end{frame}

\note{%
\textbf{\textcolor{accent}{Slide 12 : Quantification diagonale}} \hfill \textit{$\approx$ 1 min 30}

Jusqu'ici, évaluer un prédicat sur deux variables produit systématiquement
toutes les combinaisons possibles entre leurs individus. Ce comportement
par défaut n'est pourtant pas toujours celui qu'on souhaite.

Imaginons un dataset supervisé, où chaque donnée $x_i$ est associée à son
propre label $l_i$. On veut évaluer un prédicat de confiance uniquement sur
les paires qui appartiennent réellement au même exemple, $(x_0,l_0)$,
$(x_1,l_1)$, et ainsi de suite, jamais sur des combinaisons croisées comme
$(x_0,l_3)$, qui n'auraient aucun sens.

C'est exactement ce que permet la quantification diagonale : au lieu de
croiser toutes les combinaisons, elle ne conserve que les paires en
correspondance un à un, individu par individu, plutôt que toutes les
combinaisons possibles entre les deux variables. Cette correspondance
impose une condition évidente : les deux variables doivent avoir le même
nombre d'individus.

En pratique, on applique la quantification diagonale juste avant un
quantificateur. Une fois l'agrégation effectuée, le comportement normal de
croisement est restauré. Ainsi, pour tout $x, l$, $C(x,l)$ n'agrège que sur
les paires correspondantes, et cherche à répondre à la même question qu'en
classification supervisée : à quel degré le modèle attribue-t-il, en
moyenne, une bonne confiance à chaque exemple du dataset associé à son
véritable label ?
}

% ------------------------------------------------------------------
%  SLIDE 13 — Quantificateurs gardes
%  Idee unique : on peut restreindre un quantificateur a un
%  sous-ensemble d'individus via un masque booleen m(x), pas un
%  predicat flou. Exemple : un seuil de distance.
% ------------------------------------------------------------------
\begin{frame}{Quantificateurs gardés}
  \begin{center}
    \small Ils permettent de quantifier seulement sur les individus qui
    satisfont une condition \alert{booléenne} $m$, stricte.
  \end{center}

  \vspace{0.5em}
  \[
    (\forall x : m(x))\ \phi(x)
  \]
  \vspace{-0.4em}
  \[
    (\exists x : m(x))\ \phi(x)
  \]

  \vspace{0.3em}
  \begin{center}
    \small exemple : il existe une distance $d$ en dessous de laquelle
    toutes les paires sont similaires
  \end{center}
  \[
    \exists d\ \big(\forall x,y : \mathrm{dist}(x,y) < d\big)\ \mathrm{Eq}(x,y)
  \]
\end{frame}

\note{%
\textbf{\textcolor{accent}{Slide 13 : Quantificateurs gardés}} \hfill \textit{$\approx$ 1 min 15}

On souhaite parfois quantifier non pas sur l'ensemble des individus d'une
variable, mais uniquement sur ceux qui satisfont une condition
supplémentaire. Cette condition n'est pas un degré de vérité continu comme
les prédicats LTN habituels : c'est un véritable masque booléen, strictement
zéro ou un, qui sert uniquement à sélectionner qui participe à
l'agrégation.

On note $(\forall x : m(x))\,\phi(x)$ pour dire que tout $x$ satisfaisant
$m(x)$ satisfait aussi $\phi(x)$, et de même pour l'existentiel. Le masque
peut lui-même dépendre d'autres variables de la formule.

L'exemple à l'écran affirme qu'il existe une distance $d$ en dessous de
laquelle toutes les paires de points doivent être considérées comme
similaires. En pratique, cela se code avec deux arguments supplémentaires
du quantificateur : la liste des variables dont dépend la condition, et la
fonction qui calcule le masque à partir de ces variables. LTN calcule
$\mathrm{Eq}(x,y)$ normalement, calcule en parallèle ce masque, puis
n'agrège que les valeurs pour lesquelles le masque est vrai.

Cette possibilité est particulièrement utile pour concentrer le gradient
sur la partie du domaine qui vérifie effectivement la condition, plutôt que
de le diluer sur des combinaisons qui n'ont pas de sens pour la règle
logique concernée.
}

% ------------------------------------------------------------------
%  SLIDE 14 — L'apprentissage : satisfaire une base de connaissances
%  Idee unique : satisfaire K devient un probleme d'optimisation
%  continue (argmax SatAgg), sous UNE condition : chaque formule doit
%  etre close. Reutilise EstAmiDe(x,y), deja vu en slide 5.
% ------------------------------------------------------------------
\begin{frame}{Satisfaire une base de connaissances}
  \begin{center}
    \small Une base de connaissances $\mathcal{K}$ : un ensemble de
    formules qu'on veut voir vérifiées simultanément.
  \end{center}

  \vspace{0.3em}
  \[
    \mathrm{SatAgg}_{\phi \in \mathcal{K}}\, \mathcal{G}_{\theta, x \leftarrow D}(\phi)
  \]

  \vspace{0.3em}
  \begin{center}
  \footnotesize
  \renewcommand{\arraystretch}{1.4}
  \begin{tabular}{@{} >{\raggedright\arraybackslash}p{4.8cm}
                       >{\centering\arraybackslash}p{2.2cm} @{}}
    \toprule
    {\color{ash}\scshape\footnotesize formule} &
    {\color{ash}\scshape\footnotesize close ?} \\
    \midrule
    $\text{EstAmiDe}(x,y)$ & non ($x,y$ libres) \\
    $\forall x\,\text{EstAmiDe}(x,y)$ & non ($y$ libre) \\
    $\forall x\,\forall y\,\text{EstAmiDe}(x,y)$ & oui \\
    \bottomrule
  \end{tabular}
  \end{center}

  \vspace{0.2em}
  \begin{center}
    \small $\mathrm{SatAgg}$ n'agrège que des formules
    \alert{closes} : un scalaire par formule.
  \end{center}
\end{frame}

\note{%
\textbf{\textcolor{accent}{Slide 14 : L'apprentissage}} \hfill \textit{$\approx$ 2 min}

Jusqu'ici, nous avons vu comment construire et évaluer des objets LTN, mais
jamais comment faire en sorte qu'un prédicat s'améliore. C'est l'objet de
cette dernière partie : apprendre certains symboles, prédicats, fonctions,
constantes, en utilisant la satisfaction d'une base de connaissances comme
objectif d'entraînement.

Une base de connaissances est un ensemble de formules qu'on souhaite voir
vérifiées simultanément. La satisfaire ne signifie pas rendre chaque
formule strictement vraie, comme en logique classique, mais trouver les
paramètres qui rendent son degré de vérité le plus proche possible de 1.
C'est un problème d'optimisation continue, exactement comme minimiser une
perte en apprentissage automatique classique.

Notons $D$ l'ensemble complet des exemples du dataset. L'objectif s'écrit
$\mathrm{SatAgg}_{\phi \in \mathcal{K}}\ \mathcal{G}_{\theta, x \leftarrow
D}(\phi)$ : pour chaque formule $\phi$ de $\mathcal{K}$, on évalue sa
satisfaction sur les exemples $x$ de $D$ avec le modèle paramétré par
$\theta$, puis on agrège l'ensemble de ces valeurs avec $\mathrm{SatAgg}$.
Entraîner le modèle revient à choisir $\theta$ qui rend cette valeur
agrégée la plus grande possible.

Mais $\mathrm{SatAgg}$ impose une condition : chaque formule doit être
close, sans aucune variable libre, toutes capturées par un quantificateur.
Reprenons $\text{EstAmiDe}(x,y)$, déjà rencontré plus tôt : deux variables
libres, pas close. $\forall x\,\text{EstAmiDe}(x,y)$ : $y$ reste libre,
toujours pas close. $\forall x\,\forall y\,\text{EstAmiDe}(x,y)$ : les deux
sont capturées, elle est close, et son évaluation donne un unique scalaire.

C'est ce mécanisme d'agrégation, $\mathrm{SatAgg}$, que nous allons
maintenant détailler.
}

% ------------------------------------------------------------------
%  SLIDE 15 — SatAgg concretement : la perte, ce qui change / reste fixe
%  Idee unique : minimiser 1-SatAgg revient a maximiser la
%  satisfaction. Callback final au compromis neuro-symbolique
%  (slide 3 et 9) : poids apprenables changent, structure fixe.
% ------------------------------------------------------------------
\begin{frame}{De la satisfaction à la perte}
  \begin{center}
    \small $\mathrm{SatAgg}$ utilise, par défaut, le même opérateur que
    $\forall$ : $\mathrm{pMeanError}$.
  \end{center}

  \vspace{0.5em}
  \[
    \mathrm{loss} = 1 - \mathrm{SatAgg}_{\phi \in \mathcal{K}}\,
    \mathcal{G}_{\theta,x\leftarrow D}(\phi)
  \]

  \vspace{0.5em}
  \begin{center}
    \small PyTorch \textbf{minimise} ; nous voulons \textbf{maximiser} la
    satisfaction. Minimiser $1-\mathrm{SatAgg}$ revient exactement à
    maximiser $\mathrm{SatAgg}$.
  \end{center}

  \vspace{0.6em}
  \begin{center}
    \small Pendant l'entraînement : les poids apprenables
    \alert{changent} ; la structure logique des formules reste
    \alert{fixe}.
  \end{center}
\end{frame}

\note{%
\textbf{\textcolor{accent}{Slide 15 : De la satisfaction à la perte}} \hfill \textit{$\approx$ 1 min 45}

$\mathrm{SatAgg}$ agrège les degrés de vérité de toutes les formules d'une
base en une seule valeur. En pratique, on utilise le même opérateur que
pour l'universel : $\mathrm{pMeanError}$, qui est d'ailleurs le choix par
défaut. Le paramètre $p$ y joue exactement le même rôle que pour les
quantificateurs.

Puisqu'on cherche à maximiser la satisfaction globale, et que les
algorithmes d'optimisation sont conçus pour minimiser une perte, on définit
la perte comme $1 - \mathrm{SatAgg}$. Minimiser cette quantité revient
exactement à maximiser la satisfaction.

À chaque itération, on calcule d'abord les degrés de vérité des formules de
la base, la phase forward, puis on ajuste les poids par rétropropagation
pour réduire la perte, la phase backward. C'est exactement le mécanisme de
descente de gradient déjà vu, appliqué cette fois à un vrai objectif
d'entraînement.

Et ce qui change pendant cet entraînement, ce sont précisément les poids
des prédicats et fonctions apprenables, et éventuellement des constantes
entraînables. La structure logique des formules, elle, reste toujours
fixe. Nous retrouvons ici, une dernière fois, la distinction posée dès la
slide sur le compromis neuro-symbolique.
}

% ------------------------------------------------------------------
%  SLIDE 16 — Trois pieges du gradient
%  Idee unique : un operateur flou valide pour INTERROGER une formule
%  ne l'est pas forcement pour APPRENDRE dessus. Trois pathologies
%  concretes, avec exemples numeriques du notebook 2b.
% ------------------------------------------------------------------
\begin{frame}{Trois pièges du gradient}
  \begin{center}
    \small Un opérateur valide pour \alert{interroger} une formule ne
    l'est pas forcément pour \alert{apprendre} dessus.
  \end{center}

  \vspace{0.3em}
  \begin{center}
  \footnotesize
  \renewcommand{\arraystretch}{1.55}
  \begin{tabular}{@{} >{\raggedright\arraybackslash}p{4.1cm}
                       >{\raggedright\arraybackslash}p{4.1cm} @{}}
    \toprule
    {\color{ash}\scshape\footnotesize problème} &
    {\color{ash}\scshape\footnotesize cause} \\
    \midrule
    le gradient qui s'annule & $\max(z,0)$ a une dérivée nulle dès que
      $z<0$ (Łukasiewicz) \\
    le gradient à passage unique & seul l'indice extrémal reçoit un
      gradient non nul (le minimum) \\
    le gradient qui explose & $z^{1/p}$ diverge près de $z=0$
      ($\mathrm{pMeanError}$, tous les $u_i=1$) \\
    \bottomrule
  \end{tabular}
  \end{center}

  \vspace{0.4em}
  \begin{center}
    \small Même la configuration produit recommandée n'échappe pas à
    ces pièges, sur ses \alert{cas limites}.
  \end{center}
\end{frame}

\note{%
\textbf{\textcolor{accent}{Slide 16 : Trois pièges du gradient}} \hfill \textit{$\approx$ 2 min}

Un opérateur flou valide a toujours un sens quand on se contente
d'interroger une formule déjà construite. Mais ce n'est plus vrai dès
qu'on veut entraîner un modèle à partir de cette formule : de nombreux
opérateurs ont des dérivées mal adaptées à la descente de gradient.

Premier problème, le gradient qui s'annule. C'est le cas de la conjonction
de Łukasiewicz, $u \land v = \max(u+v-1, 0)$. Avec $u=0{,}3$ et $v=0{,}5$,
on a $u+v-1=-0{,}2$, donc la conjonction vaut 0, et son gradient aussi :
la dérivée de $\max(z,0)$ est nulle pour tout $z$ strictement négatif, pas
seulement au point où la formule s'annule. L'optimiseur ne reçoit alors
aucun signal pour améliorer la satisfaction.

Deuxième problème, le gradient à passage unique. C'est le comportement du
minimum : sa dérivée vaut 1 pour l'indice qui réalise effectivement le
minimum, et 0 pour tous les autres. Sur un batch de huit valeurs, un seul
individu progresse à chaque étape, les sept autres restent inchangés,
quelle que soit leur qualité.

Troisième problème, le gradient qui explose. C'est le cas de
$\mathrm{pMeanError}$ quand toutes les entrées valent exactement 1 : la
somme à l'intérieur de la racine devient nulle, et la dérivée d'une
puissance fractionnaire au voisinage de zéro diverge. Paradoxalement, le
cas le meilleur possible, la satisfaction parfaite, devient numériquement
instable.

Et la configuration produit, que nous avons recommandée, n'est pas
épargnée : la t-norme produit s'annule en $u=v=0$, la t-conorme en
$u=v=1$, l'implication de Reichenbach en $u=0,v=1$, et les deux moyennes
généralisées explosent sur leurs cas limites respectifs.
}

% ------------------------------------------------------------------
%  SLIDE 17 — La configuration stable
%  Idee unique : un decalage epsilon des entrees evite mecaniquement
%  les cas limites, sans changer le resultat ailleurs. Et : ne pas
%  fixer p trop haut pendant l'entrainement.
% ------------------------------------------------------------------
\begin{frame}{La configuration stable}
  \begin{center}
    \small L'astuce : décaler très légèrement les entrées pour
    qu'elles n'atteignent jamais exactement $0$ ou $1$.
  \end{center}

  \vspace{0.5em}
  \begin{center}
    \small cas limite en $u=0$ :\quad $u' = (1-\epsilon)\,u + \epsilon$
  \end{center}
  \vspace{0.2em}
  \begin{center}
    \small cas limite en $u=1$ :\quad $u' = (1-\epsilon)\,u$
  \end{center}

  \vspace{0.4em}
  \begin{center}
    \small avec $\epsilon$ petit, par exemple $10^{-5}$
  \end{center}

  \vspace{0.4em}
  \begin{center}
    \small Autre précaution : un $p$ \alert{trop élevé} retombe dans
    le passage unique. À éviter pendant l'entraînement.
  \end{center}
\end{frame}

\note{%
\textbf{\textcolor{accent}{Slide 17 : La configuration stable}} \hfill \textit{$\approx$ 1 min 30}

Ces problèmes ne surviennent que sur des cas limites bien identifiés, et
se corrigent facilement. L'astuce : si le cas limite survient quand une
entrée vaut 0, on remplace chaque entrée $u$ par $u' = (1-\epsilon)u +
\epsilon$. Si le cas limite survient quand une entrée vaut 1, on remplace
$u$ par $u' = (1-\epsilon)u$.

L'idée est simple : en décalant très légèrement les entrées pour qu'elles
n'atteignent jamais exactement 0 ou 1, on évite mécaniquement le point
précis où la dérivée s'annule ou explose, sans changer de façon
perceptible le résultat pour toutes les valeurs qui ne sont pas déjà sur
ce cas limite. C'est cette version corrigée qu'on appelle stable.

Une dernière précaution concerne le paramètre $p$ des moyennes
généralisées. On pourrait être tenté de fixer un $p$ élevé pour obtenir
des résultats logiquement stricts. Mais en contexte d'apprentissage,
cela devient problématique : un $p$ trop élevé transforme rapidement
l'opérateur en un opérateur à passage unique, puisque pMean tend vers le
maximum et souffre alors du même défaut que le minimum. Il est donc
recommandé de ne pas fixer un $p$ trop élevé lorsqu'on entraîne un
modèle, contrairement au cas où l'on se contente d'interroger une
formule déjà construite.
}

% ------------------------------------------------------------------
%  SLIDE 18 — Exemple minimal : formule -> reseau
%  Idee unique : synthese visuelle de toute la section 2. Le meme
%  cycle vu depuis la slide 8 (grounding) jusqu'a la slide 15 (perte),
%  boucle par le gradient qui remonte ajuster les reseaux.
% ------------------------------------------------------------------
\begin{frame}{De la base de connaissances au réseau}
  \begin{center}
  \begin{tikzpicture}[x=1cm, y=1cm]

    % --- boite 1 : la base de connaissances K -------------------------
    \draw[ash, thick, rounded corners] (-4.1,3.35) rectangle (4.1,4.85);
    \node[align=center, text width=7.8cm] at (0,4.1) {
      \footnotesize\textcolor{ash}{\scshape base de connaissances
      $\mathcal{K}$}\\[2pt]
      \footnotesize $C(a,l_a) \qquad C(b,l_b)$\\[2pt]
      \footnotesize $\forall x_1,x_2,l\,\bigl(\mathrm{Sim}(x_1,x_2)
      \implies (C(x_1,l) \leftrightarrow C(x_2,l))\bigr)$
    };
    \draw[ash, -{Stealth[length=5pt]}, thick] (0,3.35) -- (0,2.9);
    \node[ash, font=\tiny, anchor=west] at (0.15,3.12)
      {grounding $\mathcal{G}_\theta$};

    % --- boite 2 : reseaux de neurones -------------------------------
    \draw[ash, thick, rounded corners] (-4.1,0.95) rectangle (4.1,2.9);
    \node[align=center, text width=7.6cm] at (0,1.925) {
      \footnotesize\textcolor{ash}{\scshape réseaux entraînables}\\[3pt]
      \footnotesize $C$ : un MLP de sortie softmax\\[3pt]
      \footnotesize $\mathrm{Sim}$ : $\exp(-\|u-v\|^2)$
    };
    \draw[ash, -{Stealth[length=5pt]}, thick] (0,0.95) -- (0,0.5);
    \node[ash, font=\tiny, anchor=west] at (0.15,0.72)
      {connecteurs et quantificateurs flous};

    % --- boite 3 : satisfaction --------------------------------------
    \draw[accent, thick, rounded corners] (-4.1,-0.4) rectangle (4.1,0.5);
    \node[align=center, text width=7.8cm] at (0,0.05) {
      \footnotesize\textcolor{accent}{\scshape degré de satisfaction}\\[1pt]
      \footnotesize un réel dans $[0,1]$, agrégé par $\mathrm{SatAgg}$
    };
    \draw[accent, -{Stealth[length=5pt]}, thick] (0,-0.4) -- (0,-0.85);
    \node[accent, font=\tiny, anchor=west] at (0.15,-0.63)
      {$\mathrm{loss} = 1 - \mathrm{SatAgg}$};

    % --- boite 4 : perte -----------------------------------------------
    \draw[accent, thick, rounded corners] (-4.1,-1.75) rectangle (4.1,-0.85);
    \node[align=center, text width=7.8cm] at (0,-1.3) {
      \footnotesize\textcolor{accent}{\scshape perte}\\[1pt]
      \footnotesize minimisée par descente de gradient
    };

    % --- boucle du gradient : contourne largement les boites a droite --
    \draw[ash, -{Stealth[length=5pt]}, thick]
      (4.1,-1.3) to[out=0, in=0, looseness=1.4]
      node[right, font=\tiny, ash, xshift=2pt] {gradient} (4.1,1.925);

  \end{tikzpicture}
  \end{center}
\end{frame}

\note{%
\textbf{\textcolor{accent}{Slide 18 : De la base de connaissances au réseau}} \hfill \textit{$\approx$ 2 min}

Rassemblons maintenant tout ce que nous avons vu dans un seul cycle, sur un
cas concret.

Prenons 19 points dans le carré $[0,4]\times[0,4]$, deux classes possibles,
et seulement deux points déjà étiquetés, $a$ pour la classe $A$, $b$ pour
la classe $B$. La base de connaissances $\mathcal{K}$ tient en trois
formules : $a$ appartient à sa classe, $b$ à la sienne, et surtout, si deux
points sont similaires, ils doivent recevoir le même verdict. Cette
dernière règle fait tout le travail : elle propage l'information de
proche en proche, comme une lumière qui passerait d'un point à ses
voisins, bien au-delà des deux seuls points étiquetés.

C'est au grounding $\mathcal{G}_\theta$ d'associer le prédicat
d'appartenance $C$ à un MLP entraînable, dont la dernière couche utilise
un softmax pour que les deux classes restent mutuellement exclusives ; et
$\mathrm{Sim}$ à une mesure de proximité géométrique entre deux points.

Les connecteurs et les quantificateurs, rendus flous par la configuration
produit, combinent ces valeurs individuelles en un seul réel : le degré de
satisfaction de la base entière.

$\mathrm{SatAgg}$ agrège ce degré, et la perte se définit comme $1$ moins
cette valeur.

Et c'est là que la boucle se referme : le gradient remonte jusqu'aux poids
de $C$, et les ajuste pour mieux satisfaire les trois formules. Après
quelques centaines d'itérations, le modèle classe correctement les 17
points jamais étiquetés, en s'appuyant uniquement sur $a$, $b$, et la
contrainte de proximité.

C'est ce cycle complet, base de connaissances, grounding, satisfaction,
perte, gradient, qui constitue un Logic Tensor Network entraîné de bout en
bout.

\smallskip
{\raggedright\textcolor{ash}{\textit{Si la question est posée :} cet
exemple complet se trouve dans le notebook
\texttt{3-knowledgebase-and-learning.ipynb}.}\par}
}

% ============================================================
\section{Passons à la pratique}

% ------------------------------------------------------------------
%  SLIDE 19 — Pourquoi cet exemple
%  Idee unique : ce cas, emprunte a DeepProbLog (provenance seulement,
%  pas comparaison), isole la vraie difficulte (combiner, pas
%  reconnaitre), donne une structure verifiable, et une difficulte
%  reglable. Pas de spoiler sur le mecanisme (gradient/latence).
% ------------------------------------------------------------------
\begin{frame}{Pourquoi cet exemple ?}
  \begin{center}
    \small Prédire la somme de deux chiffres manuscrits : un cadre
    simple pour comparer l'apprentissage neuronal et logique.
  \end{center}

  \vspace{0.6em}
  \begin{center}
  \footnotesize
  \renewcommand{\arraystretch}{1.8}
  \begin{tabular}{@{} l >{\raggedright\arraybackslash}m{5.6cm} @{}}
    \toprule
    {\color{ash}\scshape\footnotesize ce qui rend le cas exigeant} &
    {\color{ash}\scshape\footnotesize pourquoi} \\
    \midrule
    un cas reconnu dans la littérature & introduit dans l'article
      DeepProbLog, un système neuro-symbolique \\
    aucune supervision directe du classifieur & seule la somme des
      chiffres est étiquetée, pas les chiffres eux-mêmes \\
    \bottomrule
  \end{tabular}
  \end{center}
\end{frame}

\note{%
\textbf{\textcolor{accent}{Slide 19 : Pourquoi cet exemple}} \hfill \textit{$\approx$ 1 min 30}

Ce problème n'a pas été choisi au hasard. Il s'agit d'un cas bien connu de
la littérature neuro-symbolique, notamment utilisé dans l'article
DeepProbLog. Donc dans notre étude, nous comparerons un LTN à une
baseline purement supervisée. La question qui se pose alors, c'est :
\textbf{pourquoi ce problème est-il particulièrement adapté à cette
comparaison ?}

D'abord, ce problème permet d'isoler la difficulté qui nous intéresse.
Reconnaître un chiffre manuscrit est déjà une tâche bien maîtrisée par les
réseaux de neurones. Ici, l'enjeu est surtout de combiner plusieurs
informations pour produire une réponse, ce qui nous permet de mieux
observer l'apport de la logique.

Ensuite, nous connaissons la règle du problème : la réponse est la somme
des deux chiffres. Mais cette règle est utilisée différemment selon
l'approche : la baseline doit l'apprendre à partir des exemples, tandis
que le LTN peut l'exprimer explicitement sous forme de règle logique.

Enfin, la difficulté peut être facilement augmentée, en réduisant les
données ou en augmentant le nombre d'éléments à combiner. On peut ainsi
voir comment les deux approches évoluent lorsque le problème devient plus
difficile.

C'est donc un problème simple, mais suffisamment contrôlable pour faire
une comparaison claire entre les deux approches.
}

% ------------------------------------------------------------------
%  SLIDE 20 — Formuler le probleme
%  Idee unique : deux images, une seule etiquette (leur somme), le
%  predicat Digit(x,d) et la formule existentielle qui relie les deux.
%  Fidele au notebook 4 (cas a un seul chiffre).
% ------------------------------------------------------------------
\begin{frame}{Formuler le problème}
  \begin{center}
    \small Deux images de chiffres manuscrits $x,y$ en entrée, une seule
    étiquette disponible $n$ : leur somme.
  \end{center}

  \vspace{0.4em}
  \begin{center}
  \small
  \renewcommand{\arraystretch}{1.6}
  \begin{tabular}{@{} l >{\raggedright\arraybackslash}m{5.9cm} @{}}
    \toprule
    {\color{ash}\scshape\footnotesize notion} &
    {\color{ash}\scshape\footnotesize rôle} \\
    \midrule
    $\text{Digit}(x,d)$ & vraisemblance que l'image $x$ représente le
      chiffre $d$ \\
    $n$ & la somme observée, seule étiquette fournie \\
    \bottomrule
  \end{tabular}
  \end{center}

  \vspace{0.5em}
  \[
    \exists\, d_1,d_2 : d_1+d_2=n\ \ \big(\text{Digit}(x,d_1)\land
    \text{Digit}(y,d_2)\big)
  \]
\end{frame}

\note{%
\textbf{\textcolor{accent}{Slide 20 : Formuler le problème}} \hfill \textit{$\approx$ 1 min 30}

Passons maintenant à la formulation logique du problème. On dispose de deux
images de chiffres manuscrits, $x$ et $y$, et d'une seule étiquette : leur
somme $n$. Les chiffres individuels que représentent $x$ et $y$ ne sont
jamais fournis pendant l'entraînement.

On introduit un prédicat $\text{Digit}(x,d)$, qui doit renvoyer la
vraisemblance que l'image $x$ représente le chiffre $d$. C'est ce
prédicat, sous la forme d'un réseau convolutif, que nous cherchons à
apprendre.

La formule qui relie ces éléments s'écrit : il existe deux chiffres
$d_1,d_2$ dont la somme vaut $n$, tels que $x$ représente $d_1$ et $y$
représente $d_2$. Le classifieur de chiffre n'est donc jamais entraîné
directement sur sa propre tâche : sa sortie reste une information latente,
que seule la contrainte logique sur la somme permet d'exploiter.
}

% ------------------------------------------------------------------
%  SLIDE 21 — L'axiome complet : Diag et quantification gardee
%  Idee unique : reutilise diag (slide 10) et guarded (slide 12) pour
%  fermer la formule sur un batch d'exemples (x,y,n) et ne conserver
%  que les paires (d1,d2) compatibles avec n. Figure = graphe de calcul
%  du notebook 4 (identique a la figure 13 de l'article).
% ------------------------------------------------------------------
\begin{frame}{L'axiome complet}
  \footnotesize
  \vspace{-0.7em}
  \[
    \forall\,\text{Diag}(x,y,n)\ \ \Big(\exists\, d_1,d_2 : d_1+d_2=n\ \
    \big(\text{Digit}(x,d_1)\land\text{Digit}(y,d_2)\big)\Big)
  \]
  \vspace{-1em}

  \begin{center}
  \scriptsize
  \renewcommand{\arraystretch}{0.75}
  \begin{tabular}{@{} l >{\raggedright\arraybackslash}m{6.5cm} @{}}
    \toprule
    {\color{ash}\scshape\footnotesize mécanisme} &
    {\color{ash}\scshape\footnotesize rôle ici} \\
    \midrule
    $\text{Diag}(x,y,n)$ & associe le $i$-ème exemple de $x$, $y$ et $n$,
      au lieu de les combiner librement \\
    quantification gardée & ne retient, parmi les couples $(d_1,d_2)$, que
      ceux dont la somme vaut $n$ \\
    \bottomrule
  \end{tabular}
  \end{center}

  \vspace{-1.3em}
  \begin{center}
    \includegraphics[width=0.8\textwidth]{images/semisupervised-single.png}\\[-0.15em]
    {\tiny\color{ash} Graphe de calcul symbolique pour l'addition à un
    chiffre. Source : Badreddine et al., \emph{Logic Tensor Networks},
    Figure 13.}
  \end{center}
\end{frame}

\note{%
\textbf{\textcolor{accent}{Slide 21 : L'axiome complet}} \hfill \textit{$\approx$ 1 min 30}

La formule précédente doit maintenant s'appliquer à un batch entier
d'exemples, chacun composé de deux images et d'un résultat. On retrouve
ici deux mécanismes déjà vus : la quantification diagonale et la
quantification gardée.

$\text{Diag}(x,y,n)$ garantit que le $i$-ème exemple de $x$, le $i$-ème de
$y$ et le $i$-ème de $n$ se correspondent : on n'évalue pas toutes les
combinaisons possibles entre images et résultats, seulement celles qui
proviennent du même exemple du jeu de données.

La quantification existentielle, elle, est gardée : parmi les couples
$(d_1,d_2)$, seuls ceux dont la somme vaut $n$ sont retenus dans
l'agrégation. C'est cette contrainte qui injecte directement l'information
symbolique dans le système.

Le schéma reprend le graphe de calcul complet, du grounding des images
jusqu'au degré de satisfaction final : le réseau convolutif produit les
deux distributions de probabilité, la conjonction les combine, le masque
de la contrainte $d_1+d_2=n$ sélectionne les couples valides, puis
l'existentiel et l'universel agrègent successivement, sur les couples puis
sur les exemples du batch.
}

% ------------------------------------------------------------------
%  SLIDE 22 — Entrainer : LTN contre baseline
%  Idee unique : meme reseau convolutif, deux facons de l'entrainer.
%  LTN maximise SatAgg (ici juste Forall, un seul axiome) ; la baseline
%  est supervisee directement sur la somme, sans jamais voir Digit.
%  Rappel bref du planning de p (deja vu section 2, pas de redite).
% ------------------------------------------------------------------
\begin{frame}{LTN contre baseline}
  \scriptsize
  \begin{center}
    \footnotesize\itshape Le même réseau convolutif, mais deux façons
    différentes d'apprendre.
  \end{center}

  \begin{center}
    $D$ : tout le dataset, $B$ : un mini-batch tiré de $D$ \\[0.2em]
    $\text{Satisf} = \SatAgg_{\phi\in\KB}\,\G_{\theta,x\leftarrow D}(\phi)$,
    \quad $L = 1-\SatAgg_{\phi\in\KB}\,\G_{\theta,x\leftarrow B}(\phi)$,
    \quad $\text{CE}$ : entropie croisée
  \end{center}

  \renewcommand{\arraystretch}{1.0}
  \begin{center}
  \begin{tabular}{@{} l c c @{}}
    \toprule
    & {\color{accent}\scshape\footnotesize LTN} &
      {\color{ash}\scshape\footnotesize baseline} \\
    \midrule
    données d'entrée & \multicolumn{2}{c}{identiques} \\
    tronc convolutif & \multicolumn{2}{c}{architecture identique} \\
    taille du batch & \multicolumn{2}{c}{32} \\
    optimiseur & \multicolumn{2}{c}{Adam, $lr=0{,}001$} \\
    \midrule
    objectif & maximiser $\text{Satisf}$ &
      minimiser l'erreur sur $n$ \\
    fonction de perte & $L$ & $\text{CE}$ \\
    connaissance logique & $\exists\,d_1,d_2 : d_1+d_2=n$ & aucune \\
    \midrule
    epochs & 20 & 30 \\
    \bottomrule
  \end{tabular}
  \end{center}

  \vspace{0.3em}
  \textbf{Paramètre $p$ (LTN uniquement).} Le quantificateur $\exists$ est
  rendu progressivement plus exigeant : $p=1\rightarrow2\rightarrow4
  \rightarrow6$, par paliers de 4 epochs. On part d'une agrégation proche
  de la moyenne, puis on se rapproche progressivement du maximum.
\end{frame}

\note{%
\textbf{\textcolor{accent}{Slide 22 : LTN contre baseline}} \hfill \textit{$\approx$ 2 min}

Pour comparer équitablement les deux approches, on entraîne le même
réseau convolutif, sur les mêmes données, avec le même optimiseur, mais
de deux façons différentes.

Notons $D$ l'ensemble complet des exemples du dataset. La baseline doit
prédire directement la somme : c'est un problème de classification
classique, entraîné avec une entropie croisée. Le LTN, lui, doit
maximiser $\text{Satisf}$, la satisfaction de la base de connaissance
$\KB$ sur $D$ entier, avec $\KB$ réduite au seul axiome vu précédemment.
Ce qui revient en pratique à minimiser la fonction de perte $L$ évaluée
sur un mini-batch $B$ tiré de $D$.

La baseline est entraînée pendant 30 epochs contre 20 pour le LTN,
conformément au protocole utilisé dans le papier.

Le seul hyperparamètre propre au LTN est $p$, celui du quantificateur
existentiel, déjà rencontré. Il suit un planning progressif : $p=1$ sur
les quatre premiers epochs, puis $2$, puis $4$, puis $6$ pour le reste de
l'entraînement. Au début, le classifieur de chiffres est encore proche du
hasard : un $p$ faible répartit le gradient sur toutes les combinaisons
plausibles plutôt que de tout miser sur la meilleure hypothèse du moment,
qui pourrait n'être bonne que par hasard. Le focus se resserre ensuite, au
fur et à mesure que le classifieur devient fiable.
}

% ------------------------------------------------------------------
%  SLIDE 23 — Etendre a plusieurs chiffres
%  Idee unique : meme mecanisme (Diag + quantification gardee), mais
%  sur deux nombres a deux chiffres : 4 images, 4 variables d1..d4,
%  et la contrainte devient 10 d1 + d2 + 10 d3 + d4 = n. Prepare le
%  contraste de la slide de resultats.
% ------------------------------------------------------------------
\begin{frame}{Étendre à plusieurs chiffres}
  \begin{center}
    \small La même idée, appliquée à deux nombres à deux chiffres.
  \end{center}

  \vspace{0.4em}
  \footnotesize
  \[
    \forall\,\text{Diag}(x_1,x_2,y_1,y_2,n)\ \ \Big(\exists\,
    d_1,d_2,d_3,d_4 : 10d_1+d_2+10d_3+d_4=n
  \]
  \[
    \big(\text{Digit}(x_1,d_1)\land\text{Digit}(x_2,d_2)\land
    \text{Digit}(y_1,d_3)\land\text{Digit}(y_2,d_4)\big)\Big)
  \]

  \vspace{0.3em}
  \begin{center}
  \small
  \renewcommand{\arraystretch}{1.5}
  \begin{tabular}{@{} l >{\raggedright\arraybackslash}m{4.3cm}
                       >{\raggedright\arraybackslash}m{4.3cm} @{}}
    \toprule
    & {\color{ash}\scshape\footnotesize addition à un chiffre} &
      {\color{accent}\scshape\footnotesize addition à deux chiffres} \\
    \midrule
    images par exemple & 2 (celle de $x$ et celle de $y$) & 4 (celles de
      $x_1$, $x_2$, $y_1$ et $y_2$) \\
    variables existentielles & $d_1,d_2$ & $d_1,d_2,d_3,d_4$ \\
    contrainte & $d_1+d_2=n$ & $10d_1+d_2+10d_3+d_4=n$ \\
    \bottomrule
  \end{tabular}
  \end{center}
\end{frame}

\note{%
\textbf{\textcolor{accent}{Slide 23 : Étendre à plusieurs chiffres}} \hfill \textit{$\approx$ 1 min 30}

L'expérience s'étend naturellement à des nombres comportant plusieurs
chiffres. On considère ici deux nombres à deux chiffres, donc quatre
images en entrée, $x_1,x_2$ pour le premier nombre, $y_1,y_2$ pour le
second.

Le mécanisme reste rigoureusement le même que pour un seul chiffre :
$\text{Diag}$ associe toujours le $i$-ème exemple de chaque variable, et
la quantification existentielle reste gardée. Seule la contrainte change,
puisqu'il faut maintenant reconstituer deux nombres à deux chiffres avant
de les additionner : $10d_1+d_2$ pour le premier nombre, $10d_3+d_4$ pour
le second, et leur somme doit égaler $n$.

Rien d'autre ne change dans la théorie. Le classifieur $\text{Digit}$
est le même réseau convolutif, appliqué quatre fois plutôt que deux ; il
n'a toujours accès à aucune étiquette individuelle, seulement à la somme
finale $n$.
}

% ------------------------------------------------------------------
%  SLIDE 24 — Resultats
%  Idee unique : sur un chiffre, LTN et baseline se valent ; sur deux
%  chiffres, la baseline s'effondre en test alors que le LTN generalise.
%  Chiffres tires des sorties reelles du notebook (derniere epoch).
% ------------------------------------------------------------------
\begin{frame}{Résultats}
  \begin{center}
    \small Accuracy en test, à la dernière epoch.
  \end{center}

  \vspace{0.5em}
  \begin{center}
  \renewcommand{\arraystretch}{1.6}
  \begin{tabular}{@{} l c c @{}}
    \toprule
    & {\color{accent}\scshape\footnotesize LTN} &
      {\color{ash}\scshape\footnotesize baseline} \\
    \midrule
    addition à un chiffre & $97\,\%$ & $97\,\%$ \\
    addition à deux chiffres & $95\,\%$ & \alert{$45\,\%$} \\
    \bottomrule
  \end{tabular}
  \end{center}

  \vspace{0.6em}
  \begin{center}
    \small Sur un seul chiffre, les deux approches se valent. Sur deux
    chiffres, la baseline \alert{s'effondre en test} alors que le LTN
    reste stable.
  \end{center}
\end{frame}

\note{%
\textbf{\textcolor{accent}{Slide 24 : Résultats}} \hfill \textit{$\approx$ 2 min}

Sur l'addition à un seul chiffre, les deux approches obtiennent une
accuracy de test très proche, autour de $97\,\%$. Sur ce cas le plus
simple, connaître la structure logique n'apporte donc pas d'avantage
mesurable : le problème est suffisamment contraint pour qu'un réseau
purement supervisé l'apprenne tout aussi bien.

L'écart apparaît sur l'addition à deux chiffres. Le LTN reste stable,
autour de $95\,\%$ de test. La baseline, elle, s'effondre à environ
$45\,\%$, alors même que son accuracy d'entraînement reste élevée : elle
mémorise, elle ne généralise pas.

La raison tient à la structure du problème. La baseline doit prédire
directement la somme parmi $199$ classes possibles pour deux nombres à
deux chiffres, avec peu d'exemples par classe. Le LTN, lui, apprend à
distinguer les $10$ chiffres possibles pour chaque entrée, puis la règle
logique combine cette information pour déterminer la somme. Il n'a donc
pas besoin d'apprendre directement les $199$ sommes possibles.
}

% ------------------------------------------------------------------
%  SLIDE 25 — Courbes d'entrainement
%  Idee unique : les memes chiffres que le tableau precedent, mais vus
%  au fil de l'entrainement. Deux figures issues du notebook.
% ------------------------------------------------------------------
\begin{frame}{Résultats en détail}
  \begin{center}
    {\scriptsize\color{ash} addition à un chiffre}\\[0.3em]
    \includegraphics[width=0.58\textwidth]{images/addition-1-chiffre.png}
  \end{center}

  \vspace{-0.3em}
  \begin{center}
    {\scriptsize\color{ash} addition à deux chiffres}\\[0.3em]
    \includegraphics[width=0.58\textwidth]{images/addition-2-chiffres.png}
  \end{center}
\end{frame}

\note{%
\textbf{\textcolor{accent}{Slide 25 : Résultats en détail}} \hfill \textit{$\approx$ 1 min 30}

Ces courbes montrent la même comparaison que le tableau précédent, mais
suivie au fil de l'entraînement.

Sur l'addition à un chiffre, les courbes de test de la baseline et du
LTN se superposent presque du début à la fin : les deux convergent vers
la même accuracy.

Sur l'addition à deux chiffres, l'écart apparaît clairement : la courbe
d'entraînement de la baseline continue de monter, mais sa courbe de test
plafonne bien plus bas, signe d'un surapprentissage. Le LTN, lui, garde
ses courbes de train et de test proches l'une de l'autre.

À droite, le niveau de satisfaction logique progresse par paliers, qui
correspondent exactement aux changements de $p$ déjà présentés : la
satisfaction augmente à chaque fois que $p$ devient plus exigeant.
}

% ============================================================
\section{Conclusion}

% ------------------------------------------------------------------
%  SLIDE 26 — Ce qu'il faut retenir
%  Idee unique : synthese des idees cles du parcours complet, pas de
%  contenu nouveau. Suit le fil : grounding, satisfaction/SatAgg,
%  diag/gardee, cas d'etude.
% ------------------------------------------------------------------
\begin{frame}{Ce qu'il faut retenir}
  \begin{itemize}
    \item Un LTN fixe la \alert{forme} d'un prédicat par une formule
      logique ; le \alert{contenu} qu'il calcule reste appris par le
      gradient.
    \item Le grounding $\mathcal{G}$ relie chaque symbole à un tenseur,
      et chaque formule à un réel dans $[0,1]$ : la logique devient
      différentiable.
    \item Apprendre revient à maximiser $\mathrm{SatAgg}$, la
      satisfaction agrégée d'une base de connaissances, exactement comme
      minimiser une perte.
    \item La quantification diagonale et la quantification gardée
      permettent d'injecter des contraintes structurelles précises,
      plutôt que d'agréger sans discernement.
    \item Sur le cas d'étude, cette structure logique permet au LTN de
      \alert{généraliser} là où une baseline purement supervisée
      \alert{mémorise}.
  \end{itemize}
\end{frame}

\note{%
\textbf{\textcolor{accent}{Slide 26 : Ce qu'il faut retenir}} \hfill \textit{$\approx$ 1 min 30}

Reprenons les idées centrales de ce parcours.

Un Logic Tensor Network fixe la forme d'un prédicat au moyen d'une
formule logique, mais laisse le contenu qu'il calcule être appris par le
gradient. C'est le compromis neuro-symbolique présenté en introduction.

Techniquement, cela repose sur le grounding $\mathcal{G}$ : chaque
symbole devient un tenseur, chaque formule devient un réel dans
$[0,1]$, ce qui rend toute la logique différentiable.

Apprendre un LTN, c'est alors maximiser $\mathrm{SatAgg}$, la
satisfaction agrégée d'une base de connaissances, un problème
d'optimisation continue en tout point comparable à minimiser une perte.

Deux mécanismes précisent cette agrégation : la quantification
diagonale, qui associe les bons éléments entre eux plutôt que de tout
combiner, et la quantification gardée, qui restreint l'agrégation aux
seules combinaisons compatibles avec une contrainte.

Et sur notre cas d'étude, cette structure logique a un effet concret et
mesuré : elle permet au LTN de généraliser sur l'addition à plusieurs
chiffres, là où une baseline purement supervisée se contente de
mémoriser les exemples d'entraînement.
}

% ------------------------------------------------------------------
%  SLIDE 27 — References
%  Idee unique : sources citees ou implicitement utilisees dans le
%  deck (papier LTN, papier DeepProbLog, notebooks du depot).
% ------------------------------------------------------------------
\begin{frame}{Références}
  \small
  \begin{thebibliography}{99}
    \bibitem{ltn} S.~Badreddine, A.~d'Avila Garcez, L.~Serafini,
      M.~Spranger. \emph{Logic Tensor Networks}. Artificial
      Intelligence, 2022.
    \bibitem{deepproblog} R.~Manhaeve, S.~Dumančić, A.~Kimmig,
      T.~Demeester, L.~De Raedt. \emph{DeepProbLog: Neural Probabilistic
      Logic Programming}. NeurIPS, 2018.
    \bibitem{mnist} Y.~LeCun, C.~Cortes, C.~Burges. \emph{The MNIST
      Database of Handwritten Digits}.
    \bibitem{ltntorch} LTNtorch, l'implémentation PyTorch de Logic
      Tensor Networks, et ses notebooks tutoriels.
    \bibitem{repo} Code et implémentation de cette présentation :
      \url{https://github.com/vleonel-junior/logic-tensor-networks}
  \end{thebibliography}
\end{frame}

\note{%
\textbf{\textcolor{accent}{Slide 27 : Références}} \hfill \textit{$\approx$ 30 s}

Pour finir, les sources principales de cette présentation : l'article
Logic Tensor Networks, qui fonde tout le formalisme présenté ; l'article
DeepProbLog, qui a introduit le cas d'étude de l'addition de chiffres
manuscrits ; le jeu de données MNIST, utilisé pour cette tâche ; et les
notebooks tutoriels de LTNtorch, dont sont issus les exemples et le code
présentés. Le code de cette présentation, notamment l'implémentation du
cas d'étude, est disponible sur le dépôt GitHub indiqué en dernière
ligne.
}

\end{document}
